%% CVPR paper on Derive, in the shape of the CVPR author kit
%% (https://github.com/cvpr-org/author-kit).
%%
%% cvpr.sty and ieeenat_fullname.bst come from the author kit. Derive adds them to this
%% bundle when the paper is created from the template; if they are missing here (README.md
%% says why), download them from the kit and put them next to this file.
%%
%% derive.sty is Derive's package for tables and figures that update without a new
%% version: \derivetable{name} and \derivefigure{name}. The rendered page shows each
%% slot's current value, and "Download LaTeX source" writes the current data next to this
%% file as derive-dynamic/<name>.tex, so the paper compiles with what the page shows.
\documentclass[10pt,twocolumn,letterpaper]{article}

%%%%%%%%% PAPER TYPE - update for the final version
% \usepackage{cvpr}              % camera-ready
\usepackage[review]{cvpr}      % review (anonymous, with line numbers and the paper id)
% \usepackage[pagenumbers]{cvpr} % page numbers, e.g. for an arXiv version

\usepackage{graphicx}
\usepackage{amsmath,amssymb}
\usepackage{booktabs}
\usepackage{derive}

%% hyperref last, as the kit recommends; cleveref after it.
\definecolor{cvprblue}{rgb}{0.21,0.49,0.74}
\usepackage[pagebackref,breaklinks,colorlinks,allcolors=cvprblue]{hyperref}
\usepackage[capitalize]{cleveref}
\crefname{section}{Sec.}{Secs.}
\Crefname{section}{Section}{Sections}
\crefname{table}{Tab.}{Tabs.}
\Crefname{table}{Table}{Tables}

%%%%%%%%% PAPER ID - update before submitting
\def\paperID{*****}
\def\confName{CVPR}
\def\confYear{2026}

\title{Paper Title}

\author{First Author\\
Institution1\\
Institution1 address\\
{\tt\small firstauthor@i1.org}
\and
Second Author\\
Institution2\\
First line of institution2 address\\
{\tt\small secondauthor@i2.org}
}

\begin{document}
\maketitle

\begin{abstract}
One paragraph: the problem, why it matters, what this paper does, and what the reader
gains. Replace this text.
\end{abstract}

\section{Introduction}
\label{sec:intro}

State the problem and why existing work leaves it open~\cite{example2020method}. Say
what this paper contributes, in the order the sections deliver it; \cref{fig:teaser}
should already make the point.

\begin{figure}[t]
  \centering
  \derivefigure[width=\linewidth]{teaser}
  \caption{Teaser caption: what the reader sees at a glance. The image is a dynamic
  figure: replace it from Derive's Data tab and the page and the export follow.}
  \label{fig:teaser}
\end{figure}

\section{Related Work}

Group prior work by the approach it takes rather than by year. Cite with
\verb|\cite{key}| for a numeric reference~\cite{example2023survey,example2016book}.

\section{Method}
\label{sec:method}

Introduce notation once. A displayed equation gets a number when it is referenced:
\begin{equation}
  \mathcal{L} = \sum_{i=1}^{N} \big\| f_\theta(\mathbf{x}_i) - \mathbf{y}_i \big\|_2^2 .
  \label{eq:loss}
\end{equation}
\Cref{eq:loss} is minimised over the training set; inline math such as $f_\theta$ reads
well in running text.

\section{Experiments}
\label{sec:experiments}

\Cref{tab:results} is a dynamic table: an agent (or you) fills it in through Derive's
data API as runs finish, and every reader sees the current numbers.

\begin{table}[t]
  \centering
  \derivetable{results}
  \caption{Quantitative comparison on the benchmark. Higher is better.}
  \label{tab:results}
\end{table}

A static table works too, with booktabs rules:

\begin{table}[t]
  \centering
  \begin{tabular}{@{}lrr@{}}
    \toprule
    Method & mAP & FPS \\
    \midrule
    Baseline & 41.2 & 18 \\
    Ours & \textbf{45.7} & \textbf{31} \\
    \bottomrule
  \end{tabular}
  \caption{A static table.}
  \label{tab:static}
\end{table}

\section{Conclusion}

Close with what was shown, what it does not cover, and what comes next.

{
  \small
  \bibliographystyle{ieeenat_fullname}
  \bibliography{main}
}

\end{document}
